\section{Symmetric algebraic data types}
\seclabel{algebraic data types}

Let us return to the tree lookup function $\lookup : \DFun{\BoundV
t}{\Tree} \Fun{\NodeIn{\BoundV t}} \naturals$ from \exampleref{node
lookup}. Consider the symmetry in the two recursive cases
\[
  \lookup \Big(\NodeV \ell n r\Big) (\Lft \:: p) \definedby \lookup \, \ell \, p
  \qquad
  \lookup \Big(\NodeV \ell n r\Big) (\Rgt \:: p) \definedby \lookup \, r \, p
\]
and define the canoniser\label{goLeft canoniser first occurrence} on paths by $\goLeft \, \Nil =
\goLeft (\Lft \:: p) = \stet$ and $\goLeft (\Rgt \::
p) = \swap$. Then, we can define $\lookup$ as follows, with dependent
types ruling out inaccessible cases:
\[
\lookup \definedby \fun{\Big(\NodeV \ell n r\Big)} \wlg\Bigg(
  \goLeft,
    \begin{cases}
      \Nil &\mapsto n \\
      \Lft \:: p' &\mapsto \lookup \, \ell \, p'
    \end{cases}
\Bigg)
\]
However, note that $\lookup$ and the core function are mutually
recursive: $\lookup$ is defined in terms of the core function, which
then recursively calls $\lookup$ on descent. This recursion is
structural, but via interdependence through a higher-order function.

Languages such as Agda and Idris have an implicit structural recursion
and a termination checker that tries to detect size changes.  This
type of implicit mutual recursion, called through the canoniser and
group action, is typically beyond the termination checker's ability to
detect a size change. Some languages provide eliminators explicitly,
such as Rocq and Lean, generating an appropriate
recursion or induction principle, which the user can use to:
(1) define the action on the data type; (2) compose the recursion
principle with \WLOG{} reasoning; and (3) develop induction principles
for equivariant reasoning. In this section we will isolate an
additional piece of data in addition to the data type, in the form of
an action of the group on the data type's \emph{constructors}. From
this structure we derive (1)--(3). Furthermore, an implementation that
naively defines $\lookup$ as above is inefficient. It flips the tree
each time the path changes the step direction. Our theory suggests
implementation techniques that avoid this inefficiency.

We proceed as follows. We use the well-established framework of \emph{initial
algebras}. After a brief recall of functors and their initial
algebras~(\subsecref{initial algebras}), we extend them to the
symmetric setting~(\subsecref{initial equivariant algebras}) through
\emph{symmetry strengths} that describe actions in terms of data type
constructors. They package folds of equivariant algebras as
equivariant traversals. Next, we show how to reuse canonisers for
traversals over the same constructors (\subsecref{wlog-traversals}),
and how symmetry strengths avoid unnecessary data type traversals by
progressively acting only on a single layer of
constructors~(\subsecref{partial evaluation}). Finally, we combine
these results with the previous section (\secref{dependent types}),
adapting data types to the
indexed setting~(\subsecref{indexed data types}).

\subsection{Initial algebras}
\subseclabel{initial algebras}

In \emph{initial algebra semantics}, each data type signature is
specified by a \emph{structure functor}~$F : \Set \to \Set$. For
example, trees over natural numbers and paths in them from
\exampleref{tree example} are given by the following two functors,
which closely correspond to their usual data type definitions:
\begin{align*}
  \mathbf{data}\, \Tree &=
    \Leaf \Or \Node \, \Tree \, \naturals \, \Tree
  &
  \mathbf{data}\, \TreePath &=
    \Nil \Or \TreeLabel \:: \TreePath \\
  \struct[\Tree] X &\definedby \terminal \amalg X \times \naturals \times X
  &
  \struct[\TreePath] X &\definedby \terminal \amalg \TreeLabel \times X
  \intertext{
    Each functor $F$ has an evident component-wise map $F : (X \to Y) \to (F X
    \to F Y)$, for example:
  }
  \fmap[\Tree]\, f\, (\injection_1\unitval) &\definedby \injection_1\unitval &
  \fmap[\TreePath]\, f\, (\injection_1\unitval) &\definedby \injection_1\unitval &
   \\
  \fmap[\Tree]\, f\, (\injection_2\triple{\ell}n{r}) &\definedby \injection_2\triple{f\,\ell}n{f\,r} &
  \fmap[\TreePath]\, f\, (\injection_2\pair{d}p) &\definedby \injection_2\pair{d}{f\,p}
\end{align*}

An \emph{algebra}~$\pair\X f$ for a functor $F$ is a set $\X$ together
with a function $f : F \X \to \X$, which describes the evaluation of
constructors. For example, the node-counting algebra is the
$\struct[\Tree]$-algebra $\pair\naturals {c : (\terminal \amalg \naturals
\times \naturals \times \naturals) \to \naturals}$, given by
\(
  c (\injection_1\unitval) \definedby 0
\) and \(
  c (\injection_2\triple{m_1}\_{m_2}) \definedby 1 + m_1 + m_2 \). A
\emph{homomorphism} of $F$-algebras $h : \pair\X f \to \pair\Y g$ is a
function $h : \X \to \Y$ such that $h \circ f = g \circ \fmap h$.

We interpret data types by \emph{initial} algebras, denoted by
$\pair{\fix F}{(\roll : F \fix F \to \fix F)}$, whose defining
universal property is that for every other $F$-algebra $\pair\X f$,
there is a unique homomorphism $\fold f : \pair{\fix F}{\roll} \to
\pair\X f$. This homomorphism interprets structurally recursive
functions. For example, the initial $\struct[\Tree]$-algebra, which we
label by $\pair{\Tree}{\roll[\Tree]}$, satisfies:
\begin{align*}
  \roll[\Tree] \injection_1\unitval = \Leaf
  \qquad
  \roll[\Tree] \injection_2\triple{\ell}{n}{r} = \Node\,{\ell}\,{n}\,{r}
\end{align*}
When readability is more important than precision, we are going
to elide $\roll$, and identify e.g. $\injection_1\unitval$ with
$\Leaf$, and $\injection_2\triple{\ell}{n}{r}$ with
$\Node\,{\ell}\,{n}\,{r}$.

For a given $F$-algebra $\pair\X f$, the universal property of the
initial algebra becomes the recurrence:
\[
  \fold f\, \Leaf = f(\injection_1\unitval)
  \qquad
  \fold f\, (\Node\,{\ell}\,{n}\,{r}) = f(\injection_2\triple{\fold f\,\ell}n{\fold f\,r})
\]
Through the algebra $\pair{\naturals}{c}$, we can compute the number of nodes in a tree $t$ by $\fold c\, t$ since
\[
  \fold c\, \Leaf = 0
  \qquad
  \fold c\, (\Node\,{\ell}\,{n}\,{r}) = 1 + \fold c\, \ell + \fold c\, r
\]

\subsection{Initial equivariant algebras}
\subseclabel{initial equivariant algebras}

To transfer our framework to the context of data types, it is not
sufficient to equip the initial $F$-algebra with a $\G$-action $(\act)
: \G \to \fix F \to \fix F$, as this does not tell us what to do for
other $F$-algebras. Instead, we take a more refined construct that
describes the interaction of the group with the data type
constructors:

\begin{definition}
  Let $\G$ be a group. A \emph{$\G$-symmetry-strong functor}~$F =
  \pair{\carrier F}{\gstrength[\G][F]}$, or a \emph{$\G$-strong
  functor} for short, is a functor $\carrier F : \Set \to \Set$
  together with a \emph{$\G$-symmetry strength}~$\gstrength[\G][F]$,
  or \emph{$\G$-strength} for short, which is a natural transformation
  \(
    \gstrength[\G][F]_{\X} : (\carrier\G \times \carrier F X) \to \carrier F(\carrier\G \times X)
  \)
  satisfying
  \begin{align*}
    \gstrength_{\X}(\unit, u) &= \carrier F(\fun x\pair-{\unit}x)\, u \\
    \gstrength_{\X}(g\mul h, u) &=
    \carrier F (\fun{\triple{g'}{h'}{x}}\pair{g'\mul h'}x)
    (\gstrength_{G\times\X}(g,
    \gstrength_{\X}(h, u)))
  \end{align*}
  for all $u \in \carrier F\X$ and for all $g,h \in \carrier\G$.
\end{definition}

Note that unlike group actions, symmetry strengths also feature the
group in the codomain, which allows us to specify how to act on
particular sub-terms.

\begin{example}
  \examplelabel{tree strengths}
  We can imagine two different $\Sym 2$-actions on trees, one that
  recursively flips the whole tree, and one that flips only the top
  node, with corresponding strengths:
  \begin{align*}
    \swap \act
    \parent{
      \begin{tree diagram}
        \& \& 3 \\
        \& 2 \& \\
        1 \& \& 4
        \treeLink1322
        \treeLink2231
        \treeLink2233
    \end{tree diagram}}
    &=
    \parent{
      \begin{tree diagram}
        3 \& \& \\
        \& 2 \& \\
        4 \& \& 1
        \treeLink1122
        \treeLink2231
        \treeLink2233
    \end{tree diagram}}
    &
    \swap {\act}'
    \parent{
      \begin{tree diagram}
        \& \& 3 \\
        \& 2 \& \\
        1 \& \& 4
        \treeLink1322
        \treeLink2231
        \treeLink2233
    \end{tree diagram}}
    &=
    \parent{
      \begin{tree diagram}
        3 \& \& \\
        \& 2 \& \\
        1 \& \& 4
        \treeLink1122
        \treeLink2231
        \treeLink2233
    \end{tree diagram}}
    \\
    \gstrength\pair\swap{\Node\, \ell\, n\, r} &\definedby
    \Node\, {\pair \swap r}\,n\,{\pair \swap \ell}
    &
    {\gstrength}'\pair\swap{\Node\, \ell\, n\, r}
    &\definedby \Node\,{\pair \stet r}\,n\,{\pair \stet\ell}
  \end{align*}
  The first strength propagates $\swap$ to the subtrees,
  while the second strength propagates $\stet$.

  We only need to define the strength for $\swap$ since the first axiom
  forces it at $\pair\stet u$ to be $\fmap(\fun x \pair\stet
  x) u$, i.e.,
  \(
    \gstrength\pair\stet{\Node\, \ell\, n\, r} \definedby
    {\gstrength}'\pair\stet{\Node\, \ell\, n\, r}
    \definedby
    \Node {\pair\stet\ell} n {\pair\stet r}
  \).
  For leaves, we define both strengths to act trivially
  \(
    \gstrength\pair g \Leaf \definedby
    {\gstrength}'\pair g \Leaf
    \definedby \Leaf
  \).
\end{example}

\begin{remark}\remarklabel{cartesian vs symmetry strengths}
  The notion of a symmetry strength is related to, but different from the
  usual Cartesian or monoidal strength. Recall that a \emph{Cartesian
  strength} for a functor $F : \Set \to \Set$ is a natural
  transformation
  \(
    \strength[F]_{\Ctx, \X}  : \Ctx \times F\X \to F(\Ctx \times \X)
  \)
  satisfying
  \(
    \strength_{\terminal,
      \X}(\unitval, u) = \fmap(\fun x\pair-{\unitval}x)u
  \)
  for all $u \in F\X$, and
  \(
    \strength_{\Ctx\times\Dtx, \X}((\ctx, \dtx), u) =
    \fmap(\fun{\triple {\ctx'}{\dtx'}x}\pair{\pair {\ctx'}{\dtx'}}x)
    \,\parent{\strength_{\Ctx, \Dtx\times\X}(
    \ctx, \strength_{\Dtx, \X}(\dtx, u))}
    \) for all $\ctx \in \Ctx$, and $\dtx \in \Dtx$.
    (Cf.~\applemmaref[D.1]{Cartesian strengths are symmetric strengths} in the Appendix
    for the diagrams.)
  %
  Every Cartesian strength is also a $\G$-strength, which leaves
  constructors as they are and propagates the same group element into
  the sub-terms. For example, for trees, it acts as:
  \[
    \strength_{\carrier\G, \X} \pair g \Leaf \definedby \Leaf
    \qquad
    \strength_{\carrier\G, \X} \pair g {\Node\,\ell\,n\,r}
    \definedby \Node\,{\pair g\ell}\,n\,{\pair g r}
  \]
  Symmetry strengths are an instance of the more general
  \emph{distributive laws} between a
  monad and a functor~\cite{beck:distributive-laws}. We use this
  perspective pervasively in our proofs (cf.~\appref[F]{cat
    dev}). Here, the monad is the \emph{free $\G$-action} monad, a
  special case of the \emph{Writer monad}, sometimes known as the \emph{complexity
    monad}~\cite{cenciarelli-moggi:a-syntactic-approach-to-modularity-in-denotational-semantics}.
  We use this monadic perspective in \subsecref{partial evaluation} and
  \subsecref{indexed data types}.
\end{remark}

We use $\G$-symmetry strengths to lift any action on $\X$ to an action
on $F\X$:
\begin{lemma}
  \lemmalabel{strength lifts actions}
  Let $\G$ be a group and $F$ a $\G$-strong functor. Then, for every
  $\G$-action $\A$, we have a $\G$-action $F\A$ given by:
  \[
  \carrier{F\A} \definedby \carrier F\,\carrier\A
  \qquad\smash{
    (\act[{F\A}]) : \carrier\G \times \carrier F\carrier\A \xto{\gstrength_{\carrier\A}}
    \carrier F(\carrier\G \times \carrier\A) \xto{\fmap (\act[\A])} \carrier F\carrier\A}
  \]
  Furthermore, if $h : \A \equito \B$ is equivariant, so is $\fmap h : F \A \equito F \B$.
\end{lemma}

\begin{example}
  The two tree symmetry strengths from \exampleref{tree strengths}
  induce the two expected actions: action induced by the first strength
  flips the node and propagates symmetry, while for the second strength, the action
  only flips the top node and does not propagate the symmetry further.
  Cartesian strengths induce the pointwise lifting of actions.
\end{example}

Symmetry-strong functors not only lift actions, but as we will see
shortly, also uniquely determine actions on their initial algebras.
Even more, the resulting algebra structure and the action are
compatible in the following sense:

\begin{definition}%
  \definitionlabel{equivariant F-algebra}
  Let $F$ be a $\G$-strong functor. An \emph{equivariant $F$-algebra}
  $\pair \A f$ consists of a $\G$-action $\A$ and an equivariant map
  $f : F \A \equito \A$, i.e., an $\carrier F$-algebra $\pair{\carrier\A}f$
  such that:
      \diagram{22}
  An \emph{equivariant algebra homomorphism} $h : \pair\A f \equito \pair \B
  g$ is simultaneously an $\carrier F$-homomorphism $h : \pair{\carrier\A}f \to
  \pair{\carrier\B}g$ that is also equivariant $h : \A \equito \B$.
\end{definition}

Recall that initial algebra semantics supports a modular account of
data types~\cite{swierstra:data-types-a-la-carte}: coproducts $F_1 \amalg
F_2$ of functors are given pointwise by $(F_1 \amalg F_2) X = F_1 X \amalg F_2
X$, and an algebra for this coproduct functor is given by an algebra
structure for each functor over the same carrier set. This modularity
extends to symmetry-strong functors and their equivariant algebras:
%
\begin{proposition}[modularity]\propositionlabel{modularity}
  Let $F_1, F_2 : \Set \to \Set$ be $\G$-strong functors with symmetry strengths
  $\gstrength_1, \gstrength_2$. The cotupled map
  \(
    \coseq{\strength_1, \strength_2}\pair g{\injection_ix}
    \definedby
    \injection_i\strength_i\pair g{x}
  \)
  is a $\G$-symmetry strength for the coproduct functor $\carrier F_1 \amalg \carrier F_2 : \Set \to \Set$.
  Moreover, given an action $\A$, an equivariant algebra $f : (F_1
  \amalg F_2) \A \equito \A$ for this strength amounts to a pair of
  equivariant algebra structures $f_i : F_i\A \equito \A$.
\end{proposition}

Using (cf.~\appref[D]{algebraic data types}) an
appropriate notion of parameterised initiality originally proposed by
\citeauthor{bird-paterson:generalised-folds-for-nested-datatypes}~\cite{%
  bird-paterson:generalised-folds-for-nested-datatypes},
the symmetry strength lets us equip the initial algebra with an action
that makes it the \emph{initial equivariant algebra}. Then, to define
an equivariant homomorphism $h : \fix F \equito \A$, we need to define
an equivariant $F$-algebra $g : F\A \equito \A$, possibly with
the \WLOG construction.

\begin{theorem}[equivariant initiality]
  \theoremlabel{representation for algebras}
  Let $F$ be a $\G$-strong functor and $\roll : \carrier F\fix
  \carrier F \to \fix \carrier F$ the initial algebra of $\carrier F$ in $\Set$.
  Then, there is a unique action $\fix F = \pair{\fix \carrier F}{(\act[\fix F])}$ satisfying
      \diagram{50}
  making $\pair{\fix F}\roll$ into an initial equivariant $F$-algebra, i.e., the initial $F$-algebra on $\G$-actions:
  for any equivariant $F$-algebra ${f : F\A \equito A}$, the unique
  $\carrier F$-homomorphism is equivariant: $\fold f : \fix F \equito
  \A$.
\end{theorem}

\begin{example}\examplelabel{simply typed tree path}
  \examplelabel{tree path}
  To see all the modular building blocks in action, define the path lookup function
  from \exampleref{tree example}. The mirroring symmetry of the type
  $\TreePath$ is captured by the strength:
  \begin{align*}
    \strength[\TreePath]\pair g\Nil &\definedby \Nil
    &
    \strength[\TreePath]\pair g{d \:: x} &\definedby (g \act d) \:: \pair g x
  \end{align*}
  and we can compute that the action given by
  \theoremref{representation for algebras}
  satisfies:
  \begin{align*}
    g \act (d \:: p) &=
    \big((\act) \compose (\id \times \roll)\big) \pair g{\injection_2\pair{d}{p}}
    \explain*={equivariant algebra}
    \big(\roll\compose {\struct[\TreePath](\act)} \compose \strength \big)\pair g{\injection_2\pair{d}{p}} \\&
    \explain+={strength def}
    \big(\roll\compose {\struct[\TreePath](\act)}\big)(\injection_2\pair{g \act d}{\pair g p}) =
    \roll (\injection_2\pair{g \act d}{g \act p}) =
    (g \act d) \:: g \act p
    \end{align*}
  A similar calculation for $\Nil$ shows we get the expected action
  for $\TreePath$, and a similar set of calculations shows we get the
  expected recursive action for $\Tree$ from the propagating strength
  from \exampleref{tree strengths}.
  We now construct $\lookup$ and simultaneously show it is
  equivariant as a map \(
    \TreePath \equito (\Tree \to \Maybe\naturals)
  \).
  By \theoremref{representation for
  algebras}, we can instead define an equivariant algebra:
  \[
    \lookupElim : \struct[\TreePath](\Tree \to \Maybe\naturals) \equito
    (\Tree \to \Maybe\naturals)
  \]
  Thanks to the Modularity \propositionref{modularity}, it suffices to define two
  equivariant algebra structures for each case of summand of
  $\struct[\TreePath] = \terminal \amalg \TreeLabel \times {-}$.

  The $\Nil$ summand's algebra structure
  \(
    \lookupElim_{\Nil} : \terminal \equito (\Tree \to \Maybe\naturals)
  \):
  \[
    \lookupElim_{\Nil}\unitval \, \Leaf \definedby \Nothing
    \qquad
    \lookupElim_{\Nil} \unitval \, {(\Node\, \ell\, n\, r)} \definedby \Just n
  \]
  To show that $\lookupElim_{\Nil}$ is equivariant, we refer to the fact that
  currying preserves equivariance, and calculate the behaviour of the uncurried
  version on the only interesting case:
  \[
    \lookupElim_{\Nil}\!\parent{\swap \act \pair\unitval{\Node\,\ell\,{n}\,r}}
    =
    \lookupElim_{\Nil}\!\pair\unitval{\Node\,\parent{\swap\act r}\,\!{n}\,\!\parent{\swap\act\ell}}
    =
    \Just n
    =
    \swap\act \Just n
    \]
  For $(\::)$, define the algebra
  \(
  \lookupElim_{(\::)} :
  % need this workaround to get the font right, for some reason
  \text{$\TreeLabel$} \times
    (\Tree \to \Maybe\naturals) \equito (\Tree \to
    \Maybe\naturals)
  \).
  Here, we use \WLOG reasoning, as stepping right is the same as
  stepping left on the flipped tree with the rest of the path flipped.
  Define the $\TreeLabel$-canoniser
  \(
  \canoniser' : \TreeLabel \equito \inv-{\Sym 2}
  \) that ensures canonical steps are left steps, i.e.
  \(
    \canoniser' \Lft = \stet
    \) and \(
    \canoniser' \Rgt = \swap
  \).
  Then, extend it to a canoniser~$\canoniser$ for the whole domain
  through projection via \lemmaref{projection canoniser}. Its
  canonical elements are the tuples $\seq{\Lft, f}$,
  where $f$ is the recursive call we make when traversing (necessarily left).
  Thus, we eliminated the symmetric cases in the following
  core function
  \(
  \lookLeft : \Canon\canoniser \to (\Tree \to \Maybe\naturals)
  \):
  \[
    \lookLeft\seq{\Lft, f} \definedby \begin{cases}
      \Leaf &\mapsto \Nothing \\
      \Node \, \ell \, n \, r &\mapsto  f \ell
    \end{cases}
  \]
  Now, we only need to assemble the pieces. As $\canoniser'$ is
  strictly equivariant, so is $\canoniser$ by \lemmaref{equivariance
  of projection canonisers}, and by \theoremref{all core functions
  are partially equivariant}, all core functions are equivariant,
  including $\lookLeft$. Thus $\lookupElim_{(\::)}$ is equivariant by
  \theoremref{wlog-construction}. By \propositionref{modularity}, so
  is $\lookupElim$, and by \theoremref{representation for
  algebras} so is $\lookup$.

  To see that $\lookup$ indeed matches the expected behaviour, let us
  compute its behaviour on the symmetric case, recalling that the
  exponential action conjugates, $(g \act f)(a) = g \act f(\inv g \act
  a)$:
  %
  \begin{align*}
    &\lookup \, (\Rgt \:: p) \, (\Node \, \ell \, n \, r)
    =
    \big(\fold(\lookupElim) \, (\Rgt \:: p)\big) (\Node \, \ell \, n \, r) \\
    = {}
    &\big(\lookupElim \, (\injection_2\pair\Rgt{\lookup\, p})\big) (\Node \, \ell \, n \, r)
    =
    \big(\swap \act \lookLeft\pair\Lft{\swap \act \lookup\, p}\big) (\Node \, \ell \, n \, r) \\
    = {}
    % &\swap \act \Big(\lookLeft\pair\Lft{\swap \act \lookup\, p}
    %      \big(\swap \act (\Node \, \ell \, n \, r)\big)\Big) \\
    % = {}
    &\swap \act \Big(\lookLeft\pair\Lft{\swap \act \lookup\, p}
        \big(\Node \, (\swap \act r) \, n \, (\swap \act \ell)\big)\Big) \\
    = {}
    &\swap \act \big( (\swap \act \lookup\, p)\,(\swap \act r) \big)
    =
    \swap \act \big( \swap \act (\lookup\, p \, (\swap \act \swap \act r)) \big)
    =
    \lookup \, p \, r
    \tag*{\qedhere}
  \end{align*}
\end{example}

\subsection{Data type traversals \WLOG{}}
\subseclabel{wlog-traversals}
The steps in \exampleref{simply typed tree path} are common, and we often
traverse a data type in different ways but with the same canoniser.
In this section, we present reusable abstractions that streamline these steps.

\begin{definition}
  Let $F$ be a $\G$-strong functor. A \emph{canoniser} for $F$ is a
  natural transformation $\canoniser_{\X} : \carrier F\X \to
  \carrier\G$. Each such canoniser induces the \emph{canonical
  elements functor}:
  \[
  \Canon\canoniser : \Action\G \to \Set
  \quad
  (\Canon\canoniser)\A \definedby \Canon[F\A]{\canoniser_{\carrier\A}} \definedby
  \set{
    (\canoniser_{\carrier\A}u)\act[F\A] u \suchthat u \in \carrier F\carrier\A
  }\quad
  (\Canon\canoniser)h : \fun u (\carrier Fh)u
  \]
\end{definition}

A canoniser $\canoniser$ thus induces a natural assignment of
canonisers $\canoniser_{\carrier\A} : \carrier F\carrier\A \to
\carrier \G$ for each lifted $\G$-action $F\A$.
The naturality condition means canonisers are determined by their
component at the singleton $\terminal$, i.e., by a canoniser for
$\carrier F\terminal$. The unique extension of such a canoniser
$\canoniser_{\terminal} : \carrier F\terminal \to \carrier\G$ to a
canoniser for $F$ is given by the following composition
$\canoniser_{\X} : \carrier F\X \xto{\carrier F\seq{}} \carrier
F\terminal \xto{\canoniser_{\terminal}} \carrier\G$.

\begin{example}\examplelabel{goLeft canoniser}
  The $\Sym 2$-strong functor $F_{\TreePath}$ has the
  canoniser (cf.~p.~\pageref{goLeft canoniser first occurrence})
  and canonical elements:
  \[
  \goLeft\,\Nil \definedby \stet \defines \goLeft (\Lft\:: p)
  \quad
  \goLeft (\Rgt\:: p) \definedby \swap
  \qquad
  \Canon[\A]\goLeft \definedby
  \set{\Nil, \Lft\:: a \suchthat a \in \carrier\A}
  \tag*{\qedhere}
  \]
\end{example}

We can now streamline the data we need to specify to traverse a
symmetric data type using an algebra, or to reason about such
traversals using symmetry:
\begin{definition}
  Let $F$ be a $\G$-strong functor and $\canoniser : \carrier F \to
  \carrier\G$ a canoniser for it. A \emph{$\canoniser$-core
  $F$-algebra} $\pair\A f$ consists of a $\G$-action $\A$ equipped
  with a function $f : (\Canon\canoniser)\A \to \carrier\A$. It is
  \emph{equivariant} when $f$ is equivariant $f :
  \parent{(\Canon\canoniser)\A \subset F\A} \equito \A$. A
  \emph{homomorphism} of equivariant $\canoniser$-core $F$-algebras $h
  : \pair\A f \to \pair \B g$ is an equivariant map $h : \A \equito
  \B$ satisfying: $h (f\, u) = g\, ((\Canon\canoniser) h\,u)$.
\end{definition}

We can rearrange the constructs in \exampleref{simply typed tree path}
using these concepts:
\begin{example}\examplelabel{simply typed lookup as core algebra}
  We have the following equivariant $\goLeft$-core $F_{\TreePath}$-algebra:
  \begin{gather*}
  \lookupCore : (\Canon\goLeft)(\Tree \to \Triv{\Maybe\naturals}) \to
  \carrier{\Tree \to \Triv{\Maybe\naturals}} = (\Tree\to\Maybe\naturals)
  \\
  \Nil \mapsto \begin{cases}
    \Leaf            &\mapsto \Nothing \\
    \Node \ell\,n\,r &\mapsto \Just n
  \end{cases}
  \qquad\quad
  \parent{
    \Lft\:: (\Tree\xto{\mathrm{rec} }\Maybe\naturals)
  } \mapsto \begin{cases}
    \Leaf &\mapsto \Nothing \\
    \Node\,\ell\,n\,r &\mapsto \mathrm{rec}\,\ell
  \end{cases}
  \tag*{\qedhere}
  \end{gather*}
\end{example}

These ingredients streamline the steps of \exampleref{simply typed tree path},
using the portmanteau \underline{f}old w\underline{log}:
\[
\myfigbox{$
\parent{\Canon[F\A]\canoniser_{\carrier\A} = (\Canon\canoniser)\A \xto f \carrier\A}
\mapsto
\parent{F\carrier \A \xto{\wlg(\canoniser_{\carrier\A}, f)} \carrier\A}
\mapsto
\parent{\flg[\A](\canoniser, f) :
  \fix\carrier F \xto{\fold \wlg(\canoniser_{\carrier\A}, f)} \carrier \A}
$}
\]
We chain \theoremref{wlog-construction} and the Equivariant Initiality \theoremref{representation for
  algebras}:
\begin{theorem}[core initiality]
  \theoremlabel{representation for core algebras}
  Let $F$ be a $\G$-strong functor, $\roll : \carrier F\fix
  \carrier F \to \fix \carrier F$ an initial algebra of $\carrier F$,
  and $\canoniser$ a canoniser for $F$. Then restricting $\roll$ to canonical elements
  makes $\fix F$ into the initial $\canoniser$-core $F$-algebra
  $\roll : (\Canon\canoniser)\fix F \equito \fix F$:
  for any equivariant $\canoniser$-core $F$-algebra ${f :
    (\Canon\canoniser)\A \equito A}$, the unique equivariant
  $\canoniser$-core $F$-algebra homomorphism
  into $f$ is:\/ $\flg(\canoniser, f) : \fix F \equito \A$.
\end{theorem}

This theorem implies, for example, that the $\lookup$ function from
\exampleref{simply typed lookup as core algebra} is the unique
equivariant $\goLeft$-core $F_{\TreePath}$-homomorphism given by
$\flg$.

\begin{example}\examplelabel{subtree projection}
  We reuse the canoniser to define the \emph{subtree projection} function
  \begin{gather*}
  \subtree \definedby \flg(\goLeft, \subtreeCore): \TreePath \equito (\Tree \to \Maybe\Tree)
    \\
  \subtreeCore : (\Canon\goLeft)(\Tree \to \Maybe\Tree) \equito (\Tree \to \Maybe\Tree)
  \\\!\!
  \Nil\mapsto (t \mapsto \Just\,t)
  \qquad
  (\Lft\:: \Tree\xto{\mathrm{rec}}\Maybe\Tree) \mapsto \begin{cases}
    \Leaf & \mapsto \Nothing \\
    \Node\,\ell\,n\,r & \mapsto \mathrm{rec}\,\ell
  \end{cases}
  \tag*{\qedhere}
  \end{gather*}
\end{example}

\subsection{Partially evaluating actions}
\subseclabel{partial evaluation}

Both $\lookup$ from \exampleref{simply typed tree path} and $\subtree$
from \exampleref{subtree projection} give us the result we want. But
implementations that naively follow our definitions
will repeatedly flip their arguments, and the
$\subtree$ implementation will also flip its result, leading to high time and space
complexity. Faithful to our motivation, we only demonstrate which
mathematical structures underlie a more efficient implementation, and
leave a systematic account of performance engineering to future work.

We use a general principle from modern partial evaluation folklore:
instead of performing an expensive operation repeatedly, axiomatise
its properties and design an efficient implementation for its free
algebra. Then evaluate the free algebra back into the desired
structure after partial evaluation. In our case, we shall build on top
of the \emph{free action} $\gast \X$ on a set $\X$, given by
$\carrier{\gast \X} = \carrier \G \times \X$ and $g \act[\gast \X]
\pair h x = \pair {g \mul h} x$. The free action lets us
accumulate the symmetries while leaving the input unchanged, and a
pair $\pair h x \in \carrier{\gast \X}$ captures a delayed evaluation
of $h \act[\A] x$ for some $\G$-action $\A$ with $\carrier{\A} = \X$.

In particular, consider $\gast \fix F$, the free $\G$-action on
the initial $F$-algebra. When at some point in a definition we want to
inspect $\pair g t$, we should not simply unroll $t$, as that ignores the
accumulated symmetries. We also do not want to compute $g \act t$
in implementations, since that is inefficient. What we do instead is use the symmetry
strength to act only on the first layer of the data type and
propagate the action into a delayed evaluation of a subtree
represented by the free-action monad:
\[
  \carrier{\gast \fix F}
  =
  \carrier \G \times \fix F
  \xto{\carrier\G \times \unroll}
  \carrier \G \times F \fix F
  \xto{\gstrength[F]_{\fix F}}
  F (\carrier \G \times \fix F)
  =
  F (\carrier{\gast \fix F})
\]
From \theoremref{representation for algebras}, we get $g \act t =
\fmap(\act)(\strength \pair g t)$ (ignoring unrolling), meaning that
the symmetry strength both gives us a symmetry-respecting pattern
matching of a data type, and represents a symmetry-accumulating implementation that
avoids inefficient repeated
transformations.

To see the construction in action, we can equivalently compute $\subtree \, p\, t$
with the free action:
\[
\begin{array}{@{}l@{}l@{}}
\subtree \, &p\, t =
\\&\mathrm{let}\, \pair gs = \subtreeAux\,p \pair
\stet t
\\&
\mathrel{\mathrm{in}} g\act s
\end{array}
\quad
\begin{array}{@{}l@{}}
\text{where: }\\
\quad\subtreeAux : \TreePath
\equito (\gast[\Sym 2] \Tree \to \gast[\Sym 2]\Maybe\Tree)\\
\quad \subtreeAux = \flg(\goLeft, \subtreeAuxCore)
\end{array}
\]
and $\subtreeAuxCore$
is an equivariant $\goLeft$-core
$F_{\TreePath}$-algebra over $\gast[\Sym 2] \Tree \to \gast[\Sym
2]\Maybe\Tree$:
%
\begin{gather*}
\subtreeAuxCore : \Canon\goLeft\parent{\gast[\Sym 2]\Tree \to \gast[\Sym 2]\Maybe\Tree}
\to \parent{\gast[\Sym 2]\Tree \to \gast[\Sym 2]\Maybe\Tree}\\
\begin{aligned}
\Nil &\mapsto \parent{\fun{\pair gt} \pair g{\Just t}} \\
\parent{
    \Lft\::  \parent{\begin{array}{@{}l@{}}
    \gast[\Sym 2] \Tree \xto{\mathrm{rec}}\\
    \qquad\gast[\Sym 2]\Maybe\Tree
  \end{array}}
}&\mapsto \parent{\fun{\pair gt} \begin{cases}
  \pair g \Nothing
  & (\strength \pair gt = \Leaf) \\
  \mathrm{rec}\pair{g_\ell}\ell  & (\strength \pair gt = \Node\, {\pair{g_\ell}\ell} \, n \, \pair{g_r}r)
  \end{cases}}
\end{aligned}
\end{gather*}
%
Whether the tree in the last case is flipped or not, applying the
strength ensures the correct sub-structure resides in the left
sub-tree and $g_\ell$ appropriately propagates the symmetry. More
importantly, an implementation of this design never flips the tree. For example, we have:
\begin{align*}
\subtreeAux&(\Rgt\::\Nil)\pair\stet{\Node\,\ell\,n\,r} \\&=
\swap\act\subtreeAuxCore(\Lft\::\subtreeAux (\swap \act \Nil))(\swap \act \pair\stet{\Node\,\ell\,n\,r})\\&=
\swap\act\subtreeAuxCore(\Lft\::\subtreeAux \Nil)\pair\swap{\Node\,\ell\,n\,r}\\&=
\swap\act\subtreeAux \Nil\pair\swap r=
\swap\act\pair\swap{\Just r} =
\pair{\stet}{\Just r}
\end{align*}
%
where in the reduction of $\subtreeAuxCore$, we used that $\strength
\pair\swap{\Node\,\ell\,n\,r} = \Node\,\pair\swap r\,n\,\pair\swap
\ell$. As we can see, the only evaluated actions were the free ones,
and the subtree $r$, for example, remained unflipped throughout the
argument.

There are further opportunities for optimisation. For example, instead
of stacking output symmetries $g \act -$ (in contrast to the input
ones that get absorbed by the free action) we could combine them into
a single symmetry as $g \act (h \act a) = (g \mul h) \act a$. In fact,
in this specific example, we can go even further and ignore the output
symmetries altogether. By the Core Initiality
\theoremref{representation for core algebras}, we
have $(\act[\Maybe\Tree])
\compose \subtreeAux\,p\pair gt = \subtree\,p(g \act[\Tree]t)$.
Furthermore, we can show that
$\projection_1\parent{\subtreeAux\,p\,\pair gt} = g$, hence
$\subtree\,p\,t = \projection_2\parent{\subtreeAux\,p\pair \stet t}$
and the output symmetry is never used, but is merely a reasoning
device for the correctness of the implementation.

\subsection{Indexed data types}
\subseclabel{indexed data types}

In \exampleref{simply typed tree path}, we fell back to simple types,
but thanks to the generality of the initial algebra semantics
approach, a dependently-typed development follows the same lines and
recovers the finer version from \exampleref{node lookup}.
Moreover, it uncovers nuances that degenerate in the simply-typed case. For
example, the notion of a symmetry strength must in fact be phrased in
terms of the \emph{free action} $\gast \X$ over a family. The carrier
of the free action
coincides with the Cartesian product $\carrier\G \times \X$ in the
simply-typed case, but differs in the indexed case.
As we pointed out in \remarkref{cartesian vs symmetry strengths},
viewing this product as
the free-action monad also emphasises the role of the symmetry strength
as a monad-functor distributive law.
Once we explicate
this structure, much of the development amounts to repeating the
simply-typed development with the adjective `indexed' peppered
throughout. \appref[A]{indexed adts} repeats the development, and we
describe the changes here. \figref{indexed-datatypes-summary}
summarises the difference between the simply typed theory and its
dependently typed counterpart as follows. Where appropriate, we also
describe the categorical concept.

\figfile{indexed-datatypes-summary}{Summary of symmetric data type theory}

Although the categorical theory does not capture all
concepts, e.g.~canonisers, it suffices to assign to almost every concept its
type. The theory includes indexed versions of the Modularity
\propositionref{modularity}, the Equivariant Initiality
\theoremref{representation for algebras} and Core Initiality
\theoremref{representation for core algebras}.

We generalise from the category $\Set$ to $\Fam_{\carrier\Ctx}$, the
category of $\carrier\Ctx$-indexed sets and vertical maps.  The main
technical device we use to generalise from $\Set$ to
$\Fam_{\carrier\Ctx}$ is the free $\Ctx$-indexed $\G$-action
adjunction between $\Fam_{\carrier\Ctx}$ and $\Action\G_{\Ctx}$---the
category of $\Ctx$-indexed $\G$-sets and vertical equivariant maps.
Its monadicity means we can capture its structure categorically, and
we use the existence of a right adjoint $S$ to the resulting monad $T$
over $\Fam_{\carrier\Ctx}$ in one of the proofs (cf.~\apptheoremref[F.4]{generalised parameterised initiality}).

For $\G$-action $\Ctx$, we define, for each
$\carrier\Ctx$-family $\X$, the \emph{free $\Ctx$-indexed $\G$-action over} $\X$:
\begin{gather*}
\carrier{\DPair{\mg}\G{\mg\act\X}}_{\ctx} \definedby
\DPair{\mg}{\carrier\G}{\X_{\inv \mg \act \ctx}}
\qquad
g\iact{\ctx}\pair{h}{x : \X_{\inv h\act \ctx}} \definedby
\pair{g\mul h}{x : \X_{\inv h \act \inv g \act g \act \ctx}}
\end{gather*}
We also define $\DFun{\mg}\G{\inv\mg\act\X}$,
the \emph{cofree $\Ctx$-indexed $\G$-action over} $\X$.
Its fibres are the dependent-function fibres
\(
\carrier{\DFun{\mg}\G{\inv\mg\act\X}}_{\ctx} \definedby
\DFun{\mg}{\carrier\G}{\X_{\mg \act \ctx}}
\), defined by:
\[
g\iact{\ctx}\parent{f : \DFun{\g[3]}{\carrier\G}\X_{\g[3]\act\ctx}}
\definedby\fun h f(h\mul g) \in \DFun{\mg[2]}{\carrier\G}\X_{\mg[2]\act g\act\ctx}
\]
They are the two adjoints to
the forgetful functor from $\Ctx$-indexed actions to $\Ctx$-indexed families:
\begin{theorem}\theoremlabel{descent}
  For a $\G$-action $\Ctx$ and family $\carrier\Ctx \types \X$, the
  map $\Ctx \types \fun x\pair\unit x : \X \to
  \carrier{\DPair\mg\G(\mg\act\X)}$ is the universal $\Ctx$-indexed
  $\G$-action over $\X$.  The map
  $\carrier\Ctx\types\eval\pair{-}\unit :
  \carrier{\DFun\mg\G{\inv\mg\act\X}} \to \X$ is the universal
  $\Ctx$-indexed $\G$-action under $\X$. The resulting two adjunctions
  are (co)monadic.
\end{theorem}
This theorem generalises the non-indexed case, where
$\carrier\G\times\X$ has the structure of the free $\G$-action over
$\X$.  The existence of the adjoints and their (co)monadicity follow
from abstract reasons. Let: $\oint F$ be the Grothendieck
construction for a presheaf $F : \opposite{\mathbb C} \to \Set$;
$\abs{\mathcal C}$ be the discrete category underlying a category
$\mathcal C$, and $\hat{\mathbb C}$ the presheaf category of $\mathbb
C$. Consider a $\G$-action~$\Ctx$ as a presheaf over the category with one
object and morphisms the $\G$-elements. We then have the following two equivalences. First, $\Ctx$-indexed actions
are equivalent to presheaves over $\oint\Ctx$
\(
\Action\G_{\Ctx}\simeq \widehat{\oint\Ctx}
\). Forgetting the functorial action then recovers the category of
$\carrier\Ctx$-indexed families, a presheaf category itself:
\(
\Fam_{\carrier\Ctx} \simeq \widehat{\oint{\carrier\Ctx}} = \mspace{-5mu}\widehat{\ {\abs{\oint\Ctx}}\ }
\).
The forgetful functor $\Action\G_{\Ctx} \to \Fam_{\carrier\Ctx}$
arises as precomposition by $J: {\oint\Ctx} \from \abs{{\oint\Ctx}} $.
Consequently, the two adjoints are given by (co)end
formulae, and their monadicity follows by surjectivity of this
functor~\cite[Prop.~1.1]{fiore-menni:reflective-kleisli-subcategories}.
This situation is exploited by
\citeauthor{fiore-szamozvancev:soas-theory} to account for
mathematical and computational representations of abstract syntax with
binding and substitution~\cite{fiore-szamozvancev:formal-soas,%
  fiore-szamozvancev:soas-theory,%
  szamozvancev:thesis}.
One can use the coend formula to calculate the explicit descriptions
we gave the two adjoints. We show directly
(cf.~\appref[E]{indexed data types}) that these descriptions are the
adjoints, and for completeness, we also show monadicity and comonadicity using
Beck's theorem. We then have the situation in \figref{marcelo situation}.
Therefore, $P$ is a monad, $E$ is a comonad, and $P \leftadjointto E$.
Moreover, $P^2 \definedby P\compose P \leftadjointto E^2$.
\begin{figure}
  \diagram{30}
  \caption{Categorical situation for indexed $\G$-action}\figlabel{marcelo situation}
\end{figure}
